\documentclass[10pt,a4paper]{article}
\usepackage{amsmath,amssymb,amsthm,graphicx}
% Card-local minimal preamble (journal unit sigma-2025-020-non-integrability-of-the-sasano-system-of-ty, P2-ahl16 precedent).
%% Verbatim from the pinned source: the author theorem environments and macro block,
%% so every macro and environment of the counted spans is available.
% Omitted and why: the documentclass line (the journal class sigma is not on the host;
%% the assembler supplies the class) and page-geometry/font blocks (they fight the
%% forced xelatex engine).
\usepackage{amsmath,amssymb,amsthm,mathtools}
\usepackage{xcolor}
\usepackage{url}
\usepackage[hidelinks]{hyperref}

\newcommand{\RS}[1]{\textcolor{red}{#1}}
\numberwithin{equation}{section}
\newtheorem{Theorem}{Theorem}[section]
\newtheorem*{Theorem*}{Theorem}
\newtheorem{Corollary}[Theorem]{Corollary}
\newtheorem{Lemma}[Theorem]{Lemma}
\newtheorem{Proposition}[Theorem]{Proposition}
 { \theoremstyle{definition}
\newtheorem{Definition}[Theorem]{Definition}
\newtheorem{Note}[Theorem]{Note}
\newtheorem{Example}[Theorem]{Example}
\newtheorem{Remark}[Theorem]{Remark} }
\usepackage{cleveref}
\crefname{originalreference}{original reference}{original references}
\Crefname{originalreference}{Original reference}{Original references}
\makeatletter
\newlabel{Ra}{{2.5}{source}{}{sourcepointer-1}{}}
\newlabel{Ra@cref}{{[originalreference][0][]2.5}{[1][0][]source}{}{}{}}
\makeatother
\crefname{originalreference}{original reference}{original references}
\Crefname{originalreference}{Original reference}{Original references}
\makeatletter
\newlabel{ac}{{3.5}{source}{}{sourcepointer-2}{}}
\newlabel{ac@cref}{{[originalreference][0][]3.5}{[1][0][]source}{}{}{}}
\makeatother
\crefname{originalreference}{original reference}{original references}
\Crefname{originalreference}{Original reference}{Original references}
\makeatletter
\newlabel{eq}{{3}{source}{}{sourcepointer-3}{}}
\newlabel{eq@cref}{{[originalreference][0][]3}{[1][0][]source}{}{}{}}
\makeatother
\crefname{originalreference}{original reference}{original references}
\Crefname{originalreference}{Original reference}{Original references}
\makeatletter
\newlabel{for}{{3.1}{source}{}{sourcepointer-4}{}}
\newlabel{for@cref}{{[originalreference][0][]3.1}{[1][0][]source}{}{}{}}
\makeatother
\title{The connected component of the local differential Galois group of the system}
\author{Stoyanova, Tsvetana}
\date{doi:10.3842/SIGMA.2025.020; CC BY 4.0}
\begin{document}
\maketitle
\setcounter{section}{1}
\setcounter{Theorem}{21}
\noindent\textit{Editorial note.} All delivered mathematical text is the author's own, copied from the pinned journal source file sources/sigma-2025-020-non-integrability-of-the-sasano-system-of-ty/pinned-original-source.tex (Non-Integrability of the Sasano System of Type D5(1) and Stokes Phenomena); no statement or proof step was rewritten, reordered or removed.\par\medskip
\noindent\textit{Editorial note.} The counted proof is the author's complete argument: All delivered mathematical text is the author's own, copied verbatim from the pinned SIGMA source file pinned-original-source.tex. Counted statement: Theorem environment at source lines 1145-1186 (label gal-inf), own proof at 1189-1240.\par\medskip
\noindent\textit{Editorial note.} The label gal-inf names the counted statement, and the proof environment immediately following it in the source is delivered verbatim as its proof.\par\medskip
\noindent\textit{Editorial note.} The counted statement is carried by the author's own Theorem environment in the article source.\par\medskip
\noindent\textit{Editorial note.} Journal-unit assembly: card-local preamble (author theorem environments and macros verbatim), article class, bibliography backend none; any literature citation carried by the delivered span is disclosed in the notes of this card.\par\medskip
\noindent\textit{Editorial note.} Journal numbering seed recomputed from the source itself; printed 1.22 is a cross-check only.\par\medskip
\noindent\textit{Editorial note.} Revision r002: regenerated from the same spec after the r001 assembly crashed during packaging.\par\medskip
\noindent\textit{Editorial scope.} the counted statement and its own proof environment only\par\medskip
\noindent\textit{Allowed standard prerequisites (editorial).}\begin{itemize}
\item Propositions for and ac of the same article. Supplies background that the delivered proof uses. The entry is declared as a dependency; it is not re-proved on the card, and the load-bearing steps of the counted statement are carried out inside the delivered span.
\end{itemize}\medskip
\ifdefined\hypertarget\hypertarget{sourcepointer-1}{}\fi
\noindent\textit{Source reference.} Original reference 2.5: sources/sigma-2025-020-non-integrability-of-the-sasano-system-of-ty/pinned-original-source.tex, source lines 377--377. Navigation pointer invoked at line(s) 1190: the counted text refers to this environment of the same journal unit, which is not retained in this card, so the reference is declared as a pointer carrying the number the article itself prints for it (the section and the place of the statement inside that section, derived from the shared section-scoped theorem counter of the source); the cited item is an allowed input and none of it is reproduced inside the span.\par\medskip
\ifdefined\hypertarget\hypertarget{sourcepointer-2}{}\fi
\noindent\textit{Source reference.} Original reference 3.5: sources/sigma-2025-020-non-integrability-of-the-sasano-system-of-ty/pinned-original-source.tex, source lines 956--956. Navigation pointer invoked at line(s) 1146: the counted text refers to this environment of the same journal unit, which is not retained in this card, so the reference is declared as a pointer carrying the number the article itself prints for it (the section and the place of the statement inside that section, derived from the shared section-scoped theorem counter of the source); the cited item is an allowed input and none of it is reproduced inside the span.\par\medskip
\ifdefined\hypertarget\hypertarget{sourcepointer-3}{}\fi
\noindent\textit{Source reference.} Original reference 3: sources/sigma-2025-020-non-integrability-of-the-sasano-system-of-ty/pinned-original-source.tex, source lines 640--640. Navigation pointer invoked at line(s) 1147: the counted text refers to this environment of the same journal unit, which is not retained in this card, so the reference is declared as a pointer carrying the number the article itself prints for it (the section and the place of the statement inside that section, derived from the shared section-scoped theorem counter of the source); the cited item is an allowed input and none of it is reproduced inside the span.\par\medskip
\ifdefined\hypertarget\hypertarget{sourcepointer-4}{}\fi
\noindent\textit{Source reference.} Original reference 3.1: sources/sigma-2025-020-non-integrability-of-the-sasano-system-of-ty/pinned-original-source.tex, source lines 651--651. Navigation pointer invoked at line(s) 1146: the counted text refers to this environment of the same journal unit, which is not retained in this card, so the reference is declared as a pointer carrying the number the article itself prints for it (the section and the place of the statement inside that section, derived from the shared section-scoped theorem counter of the source); the cited item is an allowed input and none of it is reproduced inside the span.\par\medskip
 \begin{Theorem}\label{gal-inf}
 With respect to the formal and actual fundamental matrix solutions, given by Propositions $\ref{for}$ and~$\ref{ac}$, the connected
 component $G^0$ of the unit element of the local differential Galois group $G$ at the origin of the system~\eqref{eq} is
 defined as follows:
 \begin{enumerate}
 \itemsep=0pt 	
 \item[$1.$]
 	If $\alpha\in\mathbb{N}$, then
 	\[
 	 G^0=\left\{\begin{pmatrix}
 	 \lambda & 0 & 0 & \mu\\
 	 0 & 1 & 0 & 0\\
 	 0 & 0 & \lambda^2 & 0\\
 	 0 & 0 & 0 & \lambda^2
 	 \end{pmatrix}, \ 	 \lambda\in\mathbb{C}^*,\, \mu\in\mathbb{C}\right\} .
 	\]
 	
 \item[$2.$]
 If $\alpha\in\mathbb{Z}_{\leq 0}$, then
 		\[
 		G^0=\left\{\begin{pmatrix}
 		\lambda & \mu & 0 & 0\\
 		0 & 1 & 0 & 0\\
 		0 & 0 & \lambda^2 & 0\\
 		0 & 0 & 0 & \lambda^2
 		\end{pmatrix}, \
 		\lambda\in\mathbb{C}^*,\, \mu\in\mathbb{C}\right\} .
 		\]
 		
 \item[$3.$]
 If $\alpha\notin\mathbb{Z}$, then
 	\[
 	G^0=\left\{\begin{pmatrix}
 	\lambda & \mu & 0 & \nu\\
 	0 & 1 & 0 & 0\\
 	0 & 0 & \lambda^2 & 0\\
 	0 & 0 & 0 & \lambda^2
 	\end{pmatrix}, \
 	\lambda\in\mathbb{C}^*,\, \mu, \nu\in\mathbb{C}\right\} .
 	\] 		
 	\end{enumerate}
 \end{Theorem}

 \begin{proof}
 From Theorem \ref{Ra}, it follows that the local differential Galois group $G$ at $\tau=0$ of the system~\eqref{eq}
 is the Zariski closure of the group generated by the formal differential Galois group
 and the Stokes matrices $St_{\pi}$ and $St_0$. Since when $\alpha\in\mathbb{Z}$ the formal monodromy $\hat{M}_0$ is equal
 to the identity matrix $I_4$ then the formal differential Galois group coincides with the exponential torus $\mathcal{T}$.
 Therefore, in this case $G$ is generated by the exponential torus and the
 Stokes matrices $St_{\pi}$ and $St_0$. Since when $\alpha\in\mathbb{N}$ the Stokes matrix $St_{\pi}$ is equal to $I_4$ the
 Galois group $G$ is generated by the exponential torus $\mathcal{T}$ and the Stokes matrix $St_0$. Hence for $\alpha\in\mathbb{N}$
 the differential Galois group coincides with its connected component $G^0$ of the unit element and
 \[
 G=G^0=\left\{\begin{pmatrix}
 \lambda & 0 & 0 & \mu\\
 0 & 1 & 0 & 0\\
 0 & 0 & \lambda^2 & 0\\
 0 & 0 & 0 & \lambda^2
 \end{pmatrix}, \ \lambda\in\mathbb{C}^*,\, \mu\in\mathbb{C}\right\} .
 \]
Similarly, when $\alpha\in\mathbb{Z}_{\leq 0}$ we have that $G$
 is generated topologically by $\mathcal{T}$ and $St_{\pi}$ and
 \[
 G=G^0=\left\{\begin{pmatrix}
 \lambda & \mu & 0 & 0\\
 0 & 1 & 0 & 0\\
 0 & 0 & \lambda^2 & 0\\
 0 & 0 & 0 & \lambda^2
 \end{pmatrix}, \ \lambda\in\mathbb{C}^*,\, \mu\in\mathbb{C}\right\} .
 \]

 When $\alpha\in\mathbb{Q}$ but $\alpha\notin\mathbb{Z}$, the formal differential Galois group is not connected since the group generated
 by $\hat{M}_0$ is not connected. However, in this case the connected component of the unit element of the formal differential
 Galois group coincides with $\mathcal{T}$. Therefore, in this case $G$ does not coincides with $G^0$ but $G^0$ is generated by $\mathcal{T}$ and
 the Stokes matrices $St_{\pi}$ and $St_0$. Hence
 \[
 G^0=\left\{\begin{pmatrix}
 \lambda & \mu & 0 & \nu\\
 0 & 1 & 0 & 0\\
 0 & 0 & \lambda^2 & 0\\
 0 & 0 & 0 & \lambda^2
 \end{pmatrix}, \ \lambda\in\mathbb{C}^*,\, \mu, \nu\in\mathbb{C}\right\} .
 \]

 In the last case when $\alpha\notin\mathbb{Q}$, we have that
 \[
 G=G^0=\left\{\begin{pmatrix}
 \lambda & \mu & 0 & \nu\\
 0 & 1 & 0 & 0\\
 0 & 0 & \lambda^2 & 0\\
 0 & 0 & 0 & \lambda^2
 \end{pmatrix}, \ \lambda\in\mathbb{C}^*,\, \mu, \nu\in\mathbb{C}\right\} .
 \]
 This ends the proof.
 \end{proof}

\end{document}
